% Generado por scripts/hoja_de_curacion.py. NO EDITAR A MANO.
\documentclass[10pt,a4paper]{article}
\usepackage[utf8]{inputenc}
\usepackage[T1]{fontenc}
\usepackage[spanish,es-nodecimaldot,es-noquoting]{babel}
\usepackage[a4paper,margin=2.1cm]{geometry}
\usepackage{lmodern}
\usepackage{booktabs}
\usepackage{array}
\usepackage{amssymb}
\usepackage{xcolor}
\usepackage{mdframed}
\usepackage{microtype}
\usepackage{titlesec}
\usepackage{fancyhdr}
\usepackage[hidelinks]{hyperref}

% LOS IDENTIFICADORES DE MATHLIB TRAEN UNICODE, y pdfLaTeX no lo sabe leer.
% El primer intento murio con «Unicode character U+2080 not set up» por un
% `min_bi` con subindice. Se declara el RANGO, no el caracter que fallo hoy:
% esta hoja se regenera, y manana los modulos pendientes seran otros.
\DeclareUnicodeCharacter{2080}{\ensuremath{_0}}
\DeclareUnicodeCharacter{2081}{\ensuremath{_1}}
\DeclareUnicodeCharacter{2082}{\ensuremath{_2}}
\DeclareUnicodeCharacter{2083}{\ensuremath{_3}}
\DeclareUnicodeCharacter{2084}{\ensuremath{_4}}
\DeclareUnicodeCharacter{2085}{\ensuremath{_5}}
\DeclareUnicodeCharacter{2086}{\ensuremath{_6}}
\DeclareUnicodeCharacter{2087}{\ensuremath{_7}}
\DeclareUnicodeCharacter{2088}{\ensuremath{_8}}
\DeclareUnicodeCharacter{2089}{\ensuremath{_9}}
\DeclareUnicodeCharacter{1D62}{\ensuremath{_i}}
\DeclareUnicodeCharacter{2C7C}{\ensuremath{_j}}
\DeclareUnicodeCharacter{207F}{\ensuremath{^n}}
\DeclareUnicodeCharacter{03B1}{\ensuremath{\alpha}}
\DeclareUnicodeCharacter{03B2}{\ensuremath{\beta}}
\DeclareUnicodeCharacter{03B3}{\ensuremath{\gamma}}
\DeclareUnicodeCharacter{03B4}{\ensuremath{\delta}}
\DeclareUnicodeCharacter{03BC}{\ensuremath{\mu}}
\DeclareUnicodeCharacter{03C3}{\ensuremath{\sigma}}
\DeclareUnicodeCharacter{03C6}{\ensuremath{\varphi}}
\DeclareUnicodeCharacter{2115}{\ensuremath{\mathbb{N}}}
\DeclareUnicodeCharacter{2124}{\ensuremath{\mathbb{Z}}}
\DeclareUnicodeCharacter{211A}{\ensuremath{\mathbb{Q}}}
\DeclareUnicodeCharacter{211D}{\ensuremath{\mathbb{R}}}
\DeclareUnicodeCharacter{2102}{\ensuremath{\mathbb{C}}}
\DeclareUnicodeCharacter{1D55C}{\ensuremath{\Bbbk}}
\DeclareUnicodeCharacter{2192}{\ensuremath{\rightarrow}}
\DeclareUnicodeCharacter{2200}{\ensuremath{\forall}}
\DeclareUnicodeCharacter{2203}{\ensuremath{\exists}}

\definecolor{acento}{RGB}{70,60,140}
\definecolor{suave}{RGB}{120,120,130}
\definecolor{fondo}{RGB}{244,243,248}
\definecolor{aviso}{RGB}{150,90,20}

\newcommand{\Mathlib}{\textsc{Mathlib}}

\titleformat{\section}{\large\bfseries\color{acento}}{}{0pt}{}
\titlespacing{\section}{0pt}{16pt}{6pt}
\titleformat{\subsection}{\normalsize\bfseries}{}{0pt}{}
\titlespacing{\subsection}{0pt}{10pt}{4pt}

\newmdenv[backgroundcolor=fondo,linewidth=0pt,skipabove=8pt,skipbelow=8pt,
          innerleftmargin=10pt,innerrightmargin=10pt,
          innertopmargin=8pt,innerbottommargin=8pt]{marca}
\newmdenv[linecolor=aviso,linewidth=1.2pt,topline=false,bottomline=false,
          rightline=false,skipabove=6pt,skipbelow=6pt,
          innerleftmargin=8pt,innertopmargin=4pt,
          innerbottommargin=4pt]{aviso}

\pagestyle{fancy}\fancyhf{}
\renewcommand{\headrulewidth}{0.3pt}
\fancyhead[L]{\small\color{suave}Metamatemático · curación pendiente}
\fancyhead[R]{\small\color{suave}\thepage}

\begin{document}
\begin{center}
{\LARGE\bfseries\color{acento} Curación pendiente}\par\smallskip
{\color{suave}\small Los 41 módulos que el grafo no cubre, con todo lo
verificable ya resuelto.}\par
{\color{suave}\footnotesize Generado por \texttt{scripts/hoja\_de\_curacion.py}
· se regenera, no se edita a mano}
\end{center}
\vspace{4pt}

\section*{No queda nada pendiente}

Los \textbf{10} módulos que la medición
sigue marcando «sin nodo» están todos decididos. Nueve son marca \texttt{T}
—ni objetos ni flechas, así que no entran— y uno ya está cubierto por un nodo
que existía. Siguen apareciendo como «sin nodo», y es correcto: no hay nodo.
Lo que no son es trabajo.\par\medskip

\noindent\small\begin{tabular}{@{}>{\raggedright\arraybackslash}p{0.34\linewidth}>{\raggedright\arraybackslash}p{0.62\linewidth}@{}}\toprule
módulo & decisión\\\midrule
\texttt{Algebraic\allowbreak{}Topology.\allowbreak{}SimplexCategory.\allowbreak{}Generators\allowbreak{}Relations.\allowbreak{}Basic} & T · es SimplexCategory presentada por generadores: la misma categoria con otro nombre, y dos nombres para un objeto gastan dos plazas\\
\texttt{Computability.\allowbreak{}AkraBazzi.\allowbreak{}SumTransform} & T · es el teorema de Akra-Bazzi y sus piezas de demostracion, no un objeto: el mismo caso que prime-factorization\\
\texttt{Control.\allowbreak{}Bitraversable.\allowbreak{}Basic} & T · interfaz de efectos de Lean: endofuntores sobre Type, y un lazo no es una dependencia\\
\texttt{Control.\allowbreak{}Fix} & T · Part.fix es el punto fijo con el que Lean define funciones parciales: maquinaria de definicion\\
\texttt{Control.\allowbreak{}Functor.\allowbreak{}Multivariate} & T · la misma rama de efectos\\
\texttt{Control.\allowbreak{}Monad.\allowbreak{}Cont} & T · la misma rama de efectos\\
\texttt{Data.\allowbreak{}TypeVec} & T · vectores de tipos para inductivos multivariados: infraestructura\\
\texttt{Order.\allowbreak{}PFilter} & T · un PFilter es un filtro sobre un preorden, pero `Order` no es area del grafo y un nodo solo no la justifica: sin padre la arista no existiria. Primer candidato si algun dia entra la teoria de ordenes\\
\texttt{SetTheory.\allowbreak{}Ordinal.\allowbreak{}Notation} & T · ONote y NONote son la forma normal de Cantor como dato computable: notacion, no objeto. El nodo `ordinals` ya existe\\
\texttt{Topology.\allowbreak{}Category.\allowbreak{}TopCat.\allowbreak{}Basic} & CUBIERTO · `TopCat` ya es la identidad de `point-set-topology`. La convencion del proyecto pone el envoltorio categorico en el campo `lean` del concepto (group-theory lleva GrpCat), no en un nodo aparte\\
\bottomrule\end{tabular}\par\medskip

Si mañana la medición destapa módulos nuevos, esta hoja vuelve a llenarse
sola.\par\medskip

\begin{marca}
\textbf{1 · La marca.} Es lo que sostiene el grafo, y no la decide nada automático:\par\smallskip
\begin{tabular}{@{}llll@{}}\toprule
\texttt{C} & una categoría & $\to$ & \textbf{vértice}\\
\texttt{S} & una subcategoría plena & $\to$ & \textbf{vértice}, y la inclusión es arista\\
\texttt{F} & un funtor o clase de flechas & $\to$ & \textbf{arista} \emph{en este ambiente}\\
\texttt{O} & un objeto individual & $\to$ & vértice degenerado\\
\texttt{T} & ni objetos ni flechas & $\to$ & \textbf{fuera}\\
\bottomrule\end{tabular}\par\smallskip
\begin{sloppypar}
\texttt{F} \textbf{no} dice «esto no puede ser un vértice nunca». Eso es falso, y está demostrado que lo es en \texttt{FlechasComoObjetos.lean}: las flechas de $C$ son exactamente los objetos de la categoría de flechas $\mathrm{Arrow}\,C$, y los funtores son los objetos de la categoría de funtores $[C,D]$. Lo que dice es que \textbf{en este grafo} —cuyos objetos son conceptos y cuyas flechas son dependencias— la etiqueta nombra una flecha.\end{sloppypar}\smallskip
\texttt{homology} es \texttt{F}: aquí no es una colección que se pueda colimitar, es el funtor \emph{a lo largo del cual} se colimita. \texttt{prime-factorization} es \texttt{T}: es un teorema, no un objeto.\par\medskip
\textbf{2 · El padre.} De qué concepto es especialización. La flecha va del general al específico.\par\medskip
\textbf{3 · Cuáles nombres se quedan.} Que exista no lo hace correcto: \texttt{Nat.Prime} es la noción de primo, \texttt{Nat.minFac} no lo es.
\end{marca}

\noindent Lo que \textbf{no} hay que decidir, porque ya está verificado: si el
nombre existe, en qué módulo vive, cuál es el canónico —las citas— y el DAG de
\texttt{import}s. Los identificadores de abajo están \textbf{leídos de su
declaración}, no deducidos de la ruta: de 447 deducidos así, 95 no existían.
\par\medskip

\begin{center}\small\begin{tabular}{@{}lr@{\quad}lr@{\quad}lr@{}}\toprule
\bottomrule\end{tabular}\end{center}
\section*{Antes de dar por buena una tanda}

\begin{marca}
\noindent\texttt{python -m scripts.recuperacion\_contra\_proofnet}\par\smallskip
Precisión y cobertura contra 371 formalizaciones de oro, con su modelo nulo y
sin gastar API.\par\medskip
\textbf{Baseline hoy: 22,8\,\% / 18,0\,\% contra 1,45\,\% / 3,3\,\% —
15,7$\times$.} Si la precisión baja, esa tanda \textbf{no entra}.\par\medskip
Ya pasó: ofrecer los sustantivos de los nodos generados bajaba de 14,0\,\% a
11,5\,\%. Añadir vocabulario tiene coste.\par\medskip
\textbf{Y la regla tiene letra pequeña.} La primera tanda de 31 nodos bajaba
la precisión a 19,9\,\% sin mover la cobertura, y la causa no era ningún
veredicto equivocado —los 47 nombres los acepta \texttt{\#check}— sino que el
emparejador tokeniza el \textbf{id y el nombre} de cada nodo:
\texttt{different-\allowbreak{}ideal} aportaba el token \texttt{different},
que sale en media biblioteca, y con eso gastaba una de las dos plazas del
prompt. La puerta que lo arregla —\emph{sólo se ocupa plaza con una keyword
declarada}— subió la línea base de 21,3\,\% a 22,8\,\%. Si una tanda baja la
precisión, mira primero si sus nodos entran en plaza por su nombre.
\end{marca}

\end{document}